\cotefile{10}
\scanpage{89}

% letterhead
INSTITUT\\
DES HAUTES ÉTUDES SCIENTIFIQUES\\
BURES-SUR-YVETTE (S.-\&-O.)\\
928-48-53

\uline{Filtration} \uline{sur foncteur fibre} \uline{pour catégorie tensorielle}

Soit \deleted{$E$} $\mathcal M$ une catégorie tensorielle \deleted{sur} ``rigide \unsure{connexe}'' sur corps $k$, $S$ un schéma sur $k$, $F$ un foncteur fibre relatif à $S$

\[
\mathcal M\longrightarrow\text{faisceaux loc libres sur } S
\]

Une filtration de $F$ par des sous-foncteurs $F^{(p)}$ est dite \uline{admissible} si elle satisfait aux conditions

(0) Chaque $F^{(p)}(M)$ est loc. fact. direct dans $M$ ($\forall M\in\mathrm{Ob}\,\mathcal M$)

\[
\text{(i) Additivité :}\quad F^{(p)}(M\ast N)=F^{(p)}(M)\ast F^{(p)}(N)
\]

\[
\text{(ii) Multiplicativité}\quad F^{(p)}(M\otimes N)=\sum_{q+r=p}F^q(M)\otimes F^r(N)
\]

(iii) $F^p(\check M)$ est l'orthogonal de $F^{-p\unsure{+1}}\deleted{\unsure{(}}(M)$

(iv) Pour chaque $M$, la filtration $F^{(p)}(M)$ est loc. exh. et sép.

Considérons alors le foncteur

\[
M\leadsto\operatorname{gr}F(M)
\]

En vertu de (0), il est à valeurs dans les modules loc libres \inserted{(gradués)} sur $S$, (i) indique qu'il est additif, (ii) qu'il est multiplicatif, et (iii) qu'il commute à $\check{\phantom M}$.

\deleted{\unsure{B}} N.B. \deleted{Ses conditions impliquent} (iii) et (iv) impliquent $F^{(0)}(1)=\underline{1}$, $F^{(1)}(1)=0$, d'où $\operatorname{gr}F(1)=\underline{1}$ réduit au degré $0$.

Supposons $\mathcal M$ $1/2$ simple, alors $\operatorname{gr}F$ \uline{exact} (sinon il faut le mettre en condition \unsure{de|en|ds|dans} (i)).

Si on pose
% archive sheet number: 88
\scanpage{90}

\mspage{\unsure{5|S|2}}

\[
G=\underline{\mathrm{Aut}}_S(F)
\]

(schéma en groupes affine sur $S$), alors $\operatorname{gr}(F)$ en tant que foncteur-fibre (en oubliant sa graduation) est défini par un $G$-torseur $P$, \inserted{par $\operatorname{gr}(F)\simeq F^P$, $P=\underline{\mathrm{Isom}}_S(F,\operatorname{gr}F)$}

De plus, la graduation sur \deleted{$\unsure{G|\otimes}$} le foncteur \deleted{sur} $\operatorname{gr}(F)$ est \unsure{défini|définie} par un hom

\[
G_{m_S}\xrightarrow{i}\deleted{\unsure{G'=}}G^P=G'.
\]

Toutes ces constructions sont can. et compatibles avec changement de base.

Parmi les isom de $F$ sur $\operatorname{gr}(F)$, on distingue ceux qui respectent la filtration \inserted{et induisent l'identité sur les gradués associés}. Ils forment (pour $S$ variable) un sous-foncteur $Q$ de $P$, qui est un \inserted{quasi-}torseur à g sous le sous-groupe $H'$ de $G'^{\deleted{P}}$ des autom qui respectent la filtration \inserted{canonique} de $F'=\operatorname{gr}F$\inserted{ et induisent l'identité sur $\operatorname{gr}\unsure{F'}$}, et un \inserted{quasi-}torseur à droite sous le sous-groupe analogue $H$ de $G$.

\margin{On suppose $\mathcal M$ de type fini}

Cependant, \inserted{pour voir que c'est un torseur,} il faut s'assurer que loc. \unsure{fppf|fpqc} $Q$ admet une section, donc qu'il existe un isom \uline{respectant les filtrations} de $F$ et $F'$ et induisant l'identité sur les gr associés.
% archive sheet number: 89
\scanpage{91}

On peut supposer déjà $F$ et $F'$ isomorphes, \deleted{\unsure{si}} s'il est vrai que les deux foncteurs fibre sont isomorphes \unsure{fppf|fpqc}. --- Il semble que ça marche du moins \deleted{en car $0$} \uline{si $k$ est de car $0$}, en supposant tjrs $\mathcal M$ semi-simple ($Q$ apparaît comme $\varprojlim$ de schémas affines et plats sur $S$, donc est affine et plat).

% slanted
\margin{Construction en termes de $G'$ \deleted{\unsure{et|\({}^{\mathrm{ad}}\)}} $G_m\xrightarrow{i}G'$ (on oublie $G$)}

Déterminons le groupe $H'$, et pour ceci choisissons un \uline{générateur} $M$, alors $G'\subset\underline{\mathrm{Aut}}_{\underline{\mathcal O}_S}(F'(M))$, et $H'$ est le sous-groupe de $G'$ formé des $g'$ tels que, $\forall p\in\mathbb Z$,

\[
g'x_{p}^{\deleted{(p)}}\equiv x_{p}^{\deleted{\unsure{(p)}}}
\deleted{\unsure{+\,g_p x_{p'}^{(p')}}}
\inserted{\bigl(F'^{(p+1)}(M)\bigr)}
\deleted{x_{p}^{\unsure{(p)}}\in}
\]

Les transformations infinitésimales de $H'$ s'obtiennent ainsi : comme $G_m$ opère sur $\mathfrak g'$ grâce à $i:G_m\longrightarrow G'$, $\mathfrak g'$ se décompose suivant les caractères de $G_m$, et on a

\deleted{$\unsure{\mathfrak g}$}

\[
\boxed{\mathfrak h'=\sum_{n\geq1}\mathfrak g'_n}
\]

% slanted
\margin{\deleted{A examiner : si $G'$ \unsure{est lisse}, $H'$ \unsure{l'est-il|l'est} également ??}}

[N.B. Le sous-groupe qui invarie la filtration sur $\unsure{D'|S'}$ correspond à l'Algèbre de Lie $\mathfrak b'=\sum_{n\geq0}\mathfrak g'_n$ --- c'est le produit semi-direct du centralisateur de $i$ dans $\mathfrak g$, et de $\mathfrak h'\ldots$].

% slanted
\margin{N.B. $H'$ est \uline{tjrs} \uline{\unsure{connexe|ouvert|connexes}} (\uline{étant} \uline{unipotent} !)}

\deleted{Le groupe} L'algèbre $\mathfrak h'$ est nilpotente, et correspond d'ailleurs à un groupe \inserted{connexe} nilpotent, tandis que $\mathfrak b'$ est parabolique ; le centralisateur de $i$ \deleted{est} est un sous-groupe de Levi.

Pb : étudier les relations entre ``ss-groupes à $1$ paramètre'' et sous-groupes paraboliques avec déc de Lévi.
% archive sheet number: 90
\scanpage{92}

% letterhead
INSTITUT\\
DES HAUTES ÉTUDES SCIENTIFIQUES\\
BURES-SUR-YVETTE (S.-\&-O.)\\
928-48-53
% archive sheet number: 91
\scanpage{93}

Bien entendu, on récupère par descente la filtration de $F$ connaissant $F'=F^P$, $i:G_m\longrightarrow G'=\underset{\underline{\mathrm{Aut}}_S(F')}{\underset{\shortparallel}{G^P}}$, et le sous-$H'$-torseur \inserted{à gauche} $Q$ de $P$.

\uuline{En résumé} (si $k$ de car. $0$, $\mathcal M$ $1/2$ simple), la donnée d'un \uline{foncteur} \inserted{\uline{fibre}} \uline{filtré} \uline{admissible} sur $\mathcal M$, à ``valeurs dans $S$'', équivaut à la donnée

(i) D'un foncteur fibre \uline{gradué} \uline{admissible} \inserted{$F'$} sur $\mathcal M$, à valeurs dans $S$, i.e. d'un foncteur fibre $F'$ sur $\mathcal M$, et d'un hom $i_1:G_{m_S}\longrightarrow G'=\underline{\mathrm{Aut}}_S F'$.

(ii) D'un $G'^{>}\deleted{\unsure{H'|G'^{\geq}}}$-torseur \deleted{$P$} $Q'$ à droite, où $\deleted{H'}G'^{>}$ \deleted{\unsure{est}} est le \deleted{\unsure{ss}} groupe uni-parabolique de $G'$ défini par $i$.

A ces données, on associe le foncteur fibre $F=F'^{Q'}$, avec la filtration déduite par descente de la filtration \deleted{canonique} de $F'$ \deleted{\unsure{gradué|qui lui}} associée \unsure{canon$^{\mathrm{t}}$} à sa graduation. D'ailleurs $Q$ est trivial sss la filtration du foncteur $F$ ``splitte'' $\ldots$

Si $\mathcal M$ est muni de\deleted{s} \inserted{la} donnée \unsure{n°} $1$ des structures \unsure{motiviques|motoriques|intrinsèques}, \deleted{\unsure{on}} alors on a $j:G_{m_S}\longrightarrow G'$ hom central, et à $i_1$ il convient de joindre $i_2$ défini par

\[
i_1(\lambda)i_2(\lambda)=j(\lambda)
\]

Alors $i_1,i_2$ sont dits \uline{complémentaires}. Les paraboliques associés sont opposés \deleted{(\unsure{$G^{>}$})}, car leur intersection est un Lévi commun.
% archive sheet number: 92
\scanpage{94}

% letterhead
INSTITUT\\
DES HAUTES ÉTUDES SCIENTIFIQUES\\
BURES-SUR-YVETTE (S.-\&-O.)\\
928-48-53
% archive sheet number: 93
\scanpage{95}

S'il y a aussi la structure motivique de Tate, définissant donc

\[
\varepsilon:G'\unsure{-|\longrightarrow}G_m
\]

alors $\varepsilon i_1(\lambda)=\lambda^n$, et l'entier $n$ est le degré du \deleted{$\unsure{1(-1)}$} ``motif de Tate ``effectif'' (\unsure{moralement|vraisblmt|réalisant|noté} $T_{\ell}(G_m)^{-1}$). Dans le cas auquel nous nous intéressons, cela doit être $1$ (on dit qu'il y a compatibilité avec la structure de Tate ---) :

\[
\varepsilon i_1(\lambda)=\lambda
\]

\uline{Exemple} Soit $S$ préschéma \inserted{connexe} sur $\mathbb Q$, \deleted{\unsure{l}} et $\mathcal M$ la catégorie des motifs sur $S$. Le \uline{foncteur de} \uline{De Rham} sur $\mathcal M$ est alors à valeurs dans les \deleted{\unsure{fais|f}}Modules loc. libres sur $S$, et est admissible.

À ce titre, il définit donc des données $F'$ (Ce n'est autre que le \uline{foncteur de Hodge}) \deleted{\unsure{défini}} $G'$, $i:G_m\longrightarrow G'$, (d'où $\unsure{G'^{>}|B'^{>}|\mathfrak g'^{>}}$) et un $G'^{>}$-torseur $Q$.

\unsure{Etant|Étant} donné la structure de $G'^{>}$ (groupe unipotent), il s'ensuit d'ailleurs que si \uline{$S$ est réduit à un pt}, $Q$ \uline{est trivial}, donc le foncteur de de-Rham filtré \uline{splitte} (mais le splitting du foncteur n'est pas canonique, en particulier n'est pas stable par les autom de ce foncteur, \unsure{ni|même|ou} par ceux qui respectent la filtration
% archive sheet number: 94
\scanpage{96}

% letterhead
INSTITUT\\
DES HAUTES ÉTUDES SCIENTIFIQUES\\
BURES-SUR-YVETTE (S.-\&-O.)\\
928-48-53

et induisent l'identité sur le gradué associé$\ldots$

Supposons maintenant $S=\operatorname{Spec} K$, $K$ un corps, et considérons un hom $i:K\longrightarrow\mathbb C$. Alors la donnée de $i$ définit le foncteur-fibre \uline{de Betti}, \deleted{soit} $F''$, à valeurs dans $\mathbb Q$.

De plus, on a des isom canoniques

\[
F''_{\mathbb C}\overset{\text{Gaga}}{\simeq} F_{\mathbb C}
\overset{\substack{\text{th. de Hodge}\\\text{des intégrales}\\\text{harmoniques}}}{\simeq}
F_{\mathbb C}^{\unsure{'|''}}
\]

D'ailleurs, l'isom $F_{\mathbb C}^{\deleted{\unsure{''}}}\simeq F'_{\mathbb C}\simeq\operatorname{gr}(F_{\mathbb C})$ est déduit de l'isom $F_{\mathbb C}\simeq F''_{\mathbb C}$, \unsure{te donnant|redonnant|traduisant} une structure \uline{réelle} sur $F_{\mathbb C}$, donnée par une involution $\sigma$, et la \uline{filtration} \inserted{can} de $F_{\mathbb C}$, par $F^{p,q}_{\mathbb C}=F^{(p)}_{\mathbb C}\cdot\sigma(F^{(q)}_{\mathbb C})$. ---
% archive sheet number: 95
\scanpage{97}

Soit $G$ un schéma en gr. affine sur corps $k$ $j:G_{m_k}-G$ un hom, i.e. une structure de $\otimes$-graduation sur le foncteur oubli $T$ sur $\operatorname{Rep}(G)$.

Cette $\otimes$-graduation définit une $\deleted{\unsure{\oplus}}\otimes$-filtration, et on peut regarder les ss-groupes

\[
U\subset P\deleted{\unsure{H}}\subset G\qquad
\deleted{\unsure{R}}P=\underline{\mathrm{Aut}}_{\otimes,\mathrm{filtr}\deleted{\unsure{,i}}}\,T,\quad
U=\underline{\mathrm{Aut}}_{\otimes,\mathrm{filt},\mathrm{id.\,sur\,Gr}}(T)
\]

Lorsque $G$ \inserted{est réductif sur $k$ de car $0$, et que} admet un module fidèle $V$, \deleted{\unsure{alors}} alors $V$ est gradué par $\mathrm{Gr}_{\deleted{\unsure{g}}j}^i(V)$, filtré par

\[
\mathrm{Filt}^{p}(V)=\sum_{i\geq p}\mathrm{Gr}_j^i(V)
\]

et $P\deleted{\unsure{H}}$ est le ss-gpe de $\underline{\mathrm{Aut}}_k(V)$ laissant \unsure{invariant|invariante} cette filtration, et $U$ le sous-gpe noyau de $P\longrightarrow\prod_i\underline{\mathrm{Aut}}_k\mathrm{Gr}_{\underline{j}}^i(V)$.

Dans le cas général, $P$ et $U$ sont des intersections de tels groupes $P_V,U_V$ ($V\in\mathrm{Ob}\,\operatorname{Rep}(\unsure{G|Q})$) ; si $P$ est de t.f. on peut prendre \uline{un} $V$ tel que $P_V=P,\ U_V=U$.

\deleted{\unsure{Prenons|Précisons}} Soit $\mathfrak g=\mathrm{Lie}\,G$, $\unsure{\xi|\zeta}\in\mathfrak g$ \deleted{\unsure{défini|définit}} l'image du générateur can de $\mathfrak t=\mathrm{Lie}\,G_m$ par $\mathrm{Lie}\,j:\deleted{\mathrm{Lie}\,G_m\to\mathfrak g}\inserted{\mathfrak t\longrightarrow\mathfrak g}$, $\mathfrak g=\sum_{\alpha\in\mathbb Z}\mathfrak g^\alpha$ la décomposition de $\mathfrak g$ suivant les car de $G_m$ opérant sur $\mathfrak g$, $V=\sum_{\beta\in\mathbb Z}V^\beta$ de $\overline{\mathrm m}$.

On a donc

\[
[\mathfrak g^\alpha,\mathfrak g^\beta]\subset\mathfrak g^{\alpha+\beta}
\qquad \mathfrak g^\alpha\cdot V^\beta\subset V^{\alpha+\beta}
\]

L'Algèbre de Lie $\mathfrak p$ de $P_V$ est formée des
% archive sheet number: 96
\scanpage{98}

$X\in\mathfrak g$ tels que

\[
X\deleted{\unsure{\mathfrak g|j}}V^\beta\subset\sum_{\beta'\geq\beta}V^{\beta'}\qquad\text{pour tt }\beta\in\mathbb Z.
\]

Il suffit que l'on ait

\[
X\in\sum_{\alpha\geq0}\mathfrak g^\alpha,
\]

et c'est aussi nécessaire si $V$ est un module fidèle, car si $X=\sum_{\alpha\in\mathbb Z}X^\alpha$, et si \deleted{$\unsure{\lambda|\alpha}$} $\mu$ est le plus petit des $\alpha$ tels que $X^\alpha\ne0$, i.e. $X=X^\mu+\sum_{\lambda>\mu}X^\lambda,\ X^\mu\ne0$,

\deleted{\unsure{si on avait}} $\exists\,\alpha$ tel. que $X^\mu\cdot V^\alpha\ne0$, \deleted{donc} et comme

\[
X^\mu V^\alpha\subset V^{\alpha+\mu}
\]

il faut qu'on ait $\unsure{v|\nu|\alpha}+\mu\leq\alpha$ i.e. $\mu\leq0$. On voit de $\overline{\mathrm m}$ que si $V$ est fidèle, on a $\mathfrak u=\mathrm{Lie}\,U=\sum_{\alpha>0}\mathfrak g^\alpha$. Par passage à la limite, on en conclut aisément, dans tous les cas

\[
\mathfrak p=\sum_{\alpha\geq0}\mathfrak g^\alpha,\qquad \mathfrak u=\sum_{\alpha>0}\mathfrak g^\alpha.
\]

\deleted{Supposons \unsure{$G$} encore que $G$ soit de t. fini, soit \unsure{un module} fidèle $V$. \unsure{Alors} On a alors}

\deleted{$V$ tel que $P=P_V,\ U=U_V$. On a donc, si $\unsure{P'|Q'}$ est le $\unsure{P'\subset P}$ \unsure{ss-gpe} parabolique de $P'=\underline{\mathrm{Aut}}(V)$, qui stabilise la filtration de $V$, et $U'$ le ss-gpe \unsure{qui induit} qui induit id sur $\mathrm{Gr}(V)$}

\[
\deleted{P=G\cap P',\qquad U=G\cap U'}
\]

\deleted{et par suite on a une imm. fermée}

\[
\deleted{G/P\unsure{\subset|\longleftarrow}G'/P'}
\]
% archive sheet number: 97
\scanpage{99}

\uline{Proposition} Supposons \deleted{$k$ parfait et $G$} \inserted{$G$ \uline{connexe} de t.f., et} réductif. Alors $P$ est un sous-groupe \uline{parabolique}, \deleted{\unsure{et}} $U$ son radical unipotent, \inserted{et leur formation commute à \unsure{tte ext. du corps|tt chgt} de base.} Dans tous les cas, $G$ \deleted{\unsure{connexe}} de t.f. $\Longrightarrow U$ unipotent$\ldots$

% slanted
\margin{\unsure{Inutile|Peut-être|Vérifier} que \uline{$G$ soit} \uline{connexe} \deleted{\unsure{pour}} \deleted{\unsure{que}}}

Supposons d'abord $G$ réductif. En car. $0$, l'assertion résulte aisément de la théorie infinitésimale. En car. $p>0$, \deleted{on \unsure{note} que l'on \unsure{est ramené|a des racines|des voisins} \unsure{de noter que} la formation de} il faut sans doute travailler un peu$\ldots$

\uline{Corollaire} Toute $\otimes$-filtration d'un foncteur fibre sur un corps parfait [ou sur un anneau quelc.] splitte $\ldots$

\deleted{\uline{\unsure{Proposition|Démonstration}}}

En effet, $H^1(k,U)=0$ par Rosenlicht$\ldots$

\inserted{\uline{Remarque}} Soit $T=\operatorname{Im} j$, et soit $U'$ un ouvert dense dans $U$. Alors (si $k$ \deleted{\unsure{parfait alg clos}} infini) l'ensemble des $uTu^{-1}$ ($u\in U'$) engendre \deleted{\unsure{le groupe} \illegible{2}} le groupe \deleted{$\unsure{P}$} $T.U$ (produit $1/2$ direct)

$\{$En effet, soit $Q$ le gpe engendré. Alors son algèbre de Lie $\underline Q$ contient les $\operatorname{ad}(u)X$ ($u\in U$), donc les \deleted{\unsure{les}} $[u,x]$, avec $u\in\mathfrak u$, donc [en car. $0$] les $\mathfrak g^\alpha,\alpha>0$ [car $[\mathfrak g^\alpha,x]=\alpha\,\mathfrak g^\alpha=\mathfrak g^\alpha$], donc $\mathfrak u\}$.

% slanted
\margin{Il est facile de donner une \unsure{dém} pas \deleted{\illegible{1}} inf. qui marche en \deleted{toute} car.}

D'autre part, pour les représentations $V$ de $G$, \deleted{\unsure{les}} on a

\[
\deleted{\unsure{V^\alpha}}\ \deleted{\unsure{V^G=}} V^{T.U}\left[=\bigcap_{u\in U'} V^{uTu^{-1}}\right]
\]

\deleted{En effet, on sait que $V^G=V^P$, $P$ étant parabolique, or $V^P=V^{T.U}$, \unsure{comme on voit en} \unsure{se}}

\deleted{en remplaçant $V$ par $V^U$ \unsure{(qui est un|(qui est une} on est ramené à voir que pour tt module $V$ sur $P/U$, on a $V^{P/U}=V^T$. \unsure{Or ceci} Or}

\deleted{si $G$ est de rang $1$, on a $\unsure{L^0}$ semi-simple de rang $1$, on a d'ailleurs}
% archive sheet number: 98
\scanpage{100}

Notes Saavedra
% archive sheet number: 99
\scanpage{111}

\deleted{$G$ schéma \inserted{en groupes, \unsure{sur un anneau|mais encore|sur une base}} \unsure{prés.|qcq} \unsure{quelconque|quasi-sép.} et plat \unsure{de type fini} sur $k$ corps}

\deleted{Tout vectoriel $V$ sur $k$ où $G$ opère est $\varinjlim$ des sous-vectoriels $V'$ de dim finie sur $k$ stables par $G$ (SGAD $\mathrm{VI}_{\mathrm B}$ 11.2.) Donc la catégorie des}

$G$ schéma en groupes \inserted{affine et \deleted{plat} libre (pas néc. fini !)} sur $S$.

\deleted{On suppose l'une des conditions suivantes satisfaites}

\deleted{a) $G$ affine}

\deleted{b) $G$ qu. cpt \unsure{quasi sép|quasi-sép.} sur $S$ et plat sur $S$}

% slanted
\margin{N.B. il \unsure{suffit|suffirait} que $A(G)$ soit ``loc'' libre pour la topologie fpqc}

\deleted{Alors} Soit $\unsure{C|\mathcal C}_{\deleted{\unsure{f|\mathrm{qc}}}}(G)$ la catégorie des $S$-modules qu. coh. où $G$ opère, $C_f(G)$ la sous-catégorie formée des $\underline E$ qui sont de t.f. Alors

\[
C(G)\overset{\approx}{\longrightarrow}\operatorname{Ind} C_f(G)
\]

cf SGAD $\mathrm{VI}_{\mathrm B}$ 11.2.

Soit alors $f:T\longrightarrow S$ un morphisme \uline{affine}, défini par $\mathscr B=f_*(\underline{\mathcal O}_T)$, \deleted{et soit $f$ un} \inserted{Nous voudrions} \deleted{A \unsure{considérer|modules}} définir la catégorie des $G_T$-\deleted{E}Modules quasi-cohérents en termes de $C(G)$ et $\mathscr B$.

Or la donnée d'un $\underline{\mathcal O}_T$-module $\underline F$ qu. coh. revient à la donnée d'un $\underline E$ \inserted{$(=f_*(\underline F))$} qu. coh. sur $S$, sur lequel $\mathscr B$ opère. On constate de plus (sans hyp. sur $G$) que la donnée d'une opération
% archive sheet number: 110
\scanpage{112}

de $G_T$ sur $\underline F$ équivaut à la donnée d'une opération de $G$ sur $\underline E$ \inserted{commutant à $\mathscr B$}, car $G_T\unsure{-|\longrightarrow}\underline{\mathrm{Aut}}(\underline F)$ équivaut à

\[
G\longrightarrow\prod_{T/S}\underline{\mathrm{Aut}}(\underline F),\quad\text{or}
\]

\[
\prod_{T/S}\underline{\mathrm{Aut}}(\underline F)\simeq
\underline{\mathrm{Aut}}_{\mathscr B}\Bigl(\prod_{T/S}\underline F\Bigr)
=\underline{\mathrm{Aut}}_{\mathscr B}(\underline E).
\]

Donc, la catégorie $C(G_T)$ équivaut à la catégorie des objets de $C(G)$ munis d'une opération de $\mathscr B$, ce qui, si $S$ est affine d'anneau $A$, donc $\unsure{B|\mathscr B|T|\mathcal B}$ défini par $A$-algèbre $B$, équivaut à la \deleted{\unsure{caté|catég.}} donnée d'un objet de $C(G)$ à anneau $B$ d'endomorphismes.

Lorsque $G$ satisfait l'une des conditions a) ou b) (p.ex. $G$ qu. cpct sur un \uline{corps} $k$), \deleted{\unsure{les|la}} catégorie $C(G_T)$ est définie en termes de $C_f(G)$.

De $\overline{\mathrm m}$, si $\varphi_G$ \inserted{$(\varphi_G^f)$} désigne le foncteur oubli de \inserted{$C(G)$} ($C_f(G)$), alors $(C(G_T),\varphi_{G_T})$ est déterminé en termes de $(C_f(G),\varphi_G^f)$.

Et il en est donc de $\overline{\mathrm m}$ de $C_f(G_T),\varphi_{G_T})$, sous-catégorie pleine de $C(G_T)$ formée des $E$ tels que $\varphi_{G_T}(E)$ de d.f. sur $B$.
% archive sheet number: 111
\scanpage{113}

Soit $G$ un schéma \uline{affine} \inserted{\unsure{loc. lib.|loc. lis.}} sur un anneau $A$.

Dans quelle mesure $G$ se récupère-t-il par la connaissance de $({C}_{\deleted{\unsure{f}}}(G),\varphi_G^{\deleted{f}})$, où $C_{\deleted{f}}(G)$ est regardé comme une $\otimes$-catégorie [pour $\otimes$ \ldots], et $\varphi_G^{\deleted{f}}$ comme un $\otimes$-foncteur.

Il faudrait récupérer $G(B)$ pour les $A$-algèbres $B$. Or pour ceci, il suffirait de récupérer $G(A)$. On a un hom. évident

\diagram{\begin{tikzcd}[ampersand replacement=\&,row sep=1.5em]
G(A)\arrow[r] \& \operatorname{Aut}_{\otimes}(\varphi_G^{\deleted{f}})\arrow[d,phantom,"{\deleted{\unsure{\wr|\shortparallel}}}" description]\\
\& \deleted{\operatorname{Aut}_{\otimes}(\varphi_G)}.
\end{tikzcd}}

On présume que c'est injectif : regarder les effets de $g\in G(A)$ \unsure{,} sur \uline{l'algèbre \inserted{affine $\unsure{\varepsilon|\mathscr E|E}$} de} $G$ !

On présume sans doute aussi que c'est surjectif, en regardant encore l'effet d'un $\otimes$-Aut $u$ de $\varphi_G^{\deleted{f}}$ sur $\varepsilon$, \deleted{ce qui donne un \unsure{autom.}} \uline{à détailler}

\uline{Question} À quelle condition une $\otimes$-catégorie \inserted{$A$-linéaire $C_{\deleted{f}}$} munie d'un $\otimes$-foncteur $\varphi^{\deleted{f}}:C_{\deleted{f}}\unsure{-|\to}\operatorname{Mod}_{f}(A)$, \deleted{\unsure{peut|pro}}\unsure{provient-il|provien-il} d'un schéma en groupes affine sur $A$ ???

Il faudrait supposer que $C$ est un Ab-topos, et que $\varphi$ commute aux $\varinjlim$

% archive sheet number: 112
